% Figure: why a fixed-step rule fails on an oscillatory integrand and
% Filon-type quadrature succeeds. Top panel: the integrand
% f(x) = cos(40 x) on [0,1] with a coarse trapezoidal sampling that
% aliases. The point is that equal-weight samples on an oscillatory
% integrand read a value governed by where the samples land, not by the
% (small) true integral.
\begin{figure}[htbp]
\centering
\begin{tikzpicture}
\begin{axis}[
  width=12cm, height=5.6cm,
  xlabel={$x$}, ylabel={$g(x)\cos(\omega x)$},
  xmin=0, xmax=1, ymin=-1.15, ymax=1.15,
  domain=0:1, samples=400,
  axis lines=middle, grid=both,
  major grid style={engblack!12}, font=\small,
  legend pos=south east, legend style={font=\footnotesize},
]
% the rapidly oscillating integrand, with a slow envelope g(x) = 1 - x/2
\addplot[engblue, thick] {(1 - x/2)*cos(deg(40*x))};
\addlegendentry{$g(x)\cos(40x)$, $g(x)=1-x/2$}
% a coarse equispaced sampling at n = 10 panels: the dots land where the
% trapezoidal rule reads, missing the oscillation entirely
\addplot[only marks, mark=*, engred, mark size=1.6pt]
  coordinates {(0,1)(0.1,0.4197)(0.2,-0.6536)(0.3,0.6997)(0.4,-0.4521)
  (0.5,0.1455)(0.6,0.2620)(0.7,-0.5346)(0.8,0.5680)(0.9,-0.3852)(1.0,0.1634)};
\addlegendentry{$n=10$ equispaced nodes}
\end{axis}
\end{tikzpicture}
\caption[Oscillatory integrand and the failure of equal-weight
quadrature]{A slowly enveloped, rapidly oscillating integrand
$g(x)\cos(\omega x)$ with $\omega = 40$. An equispaced rule with only
ten panels (red dots) samples the oscillation at scattered phases and
reads an integral governed by where the nodes happen to fall rather
than by the true value, which is small because the oscillation nearly
cancels. A polynomial-fit rule needs of order $\omega/\pi$ panels per
unit length to resolve the oscillation, so its cost grows linearly with
frequency. Filon-type quadrature integrates the oscillatory factor
$\cos(\omega x)$ analytically against a polynomial fit of the slow
envelope $g$, so its accuracy improves as $\omega$ grows rather than
degrading, the reverse of the equal-weight rule's behaviour.}
\label{fig:vol02:ch06:oscillatory-quadrature}
\end{figure}
